%! Author = lili
%! Date = 6/14/26

\newglossaryentry{kfrage}{
    name={test},
    description={todo},
    symbol={2015},
    type=frage
}

\newglossaryentry{genetKarten}{
    name={Was sind genetische Karten, nenne zwei Varianten sowie die Herleitung},
    description={todo},
    symbol={2015},
    type=frage
}

\newglossaryentry{bandenmuster}{
    name={Zu erwartende Bandenmuster von der T-DNA-Insertion in verschiedenen Pflanzen skizzieren},
    description={todo},
    symbol={2015},
    type=frage
}

\newglossaryentry{genexpressiontranskript}{
    name={Vier Methoden der Genexpressionsanalyse anhand der Transkripte},
    description={todo},
    symbol={2015},
    type=frage
}

\newglossaryentry{rekomnbinationsfrequenzen}{
    name={Rekomnbinationsfrequenzen berechnen und deuten},
    description={todo},
    symbol={2015},
    type=frage
}

\newglossaryentry{gesundbistumor}{
    name={Reihenfolge anordnen: Von gesunder Zelle bis zum metastasierenden Tumor},
    description={todo},
    symbol={2015},
    type=frage
}

\newglossaryentry{zelluläreprozesse}{
    name={Übergeordnete zelluläre Prozesse zu Moleklülen zurdnen},
    description={todo bspw. tRNA zu Translation},
    symbol={2015},
    type=frage
}

\newglossaryentry{rtpcrprimer}{
    name={Drei verschiedenen Arten von Primern für eine RT-PCR und deren Auswirkungen auf die cDNA},
    description={todo},
    symbol={2015},
    type=frage
}

\newglossaryentry{dnareplikation}{
    name={Abbildung der DNA-Replikation beschriften},
    description={todo},
    symbol={2015},
    type=frage
}

\newglossaryentry{epigenetischeModifikationen}{
    name={Multiple Choice - epigenetische Modifikationen am Histon, an welchen Untereinheiten findet welche ...-ylierung statt},
    description={todo},
    symbol={2015},
    type=frage
}

\newglossaryentry{transkriptionbistranslation}{
    name={Die vier wensentlichen Prozessierungsschritte von Transkription bis Anfang der Translation},
    description={todo},
    symbol={2015},
    type=frage
}


\newglossaryentry{wurzelbalsgalle}{
    name={Lückentext zur Wurzelbalsgalle ausfüllen},
    description={6 Wörter in Lückentext einsetzen. A. thumefaciens infiziert Pflanze: infiziert, Opiate, T-DNA, Nährstoffe, (Ti-)Plasmid, insertiert},
    symbol={2015, 2018},
    type=frage
}

\newglossaryentry{konzentrationsberechnung}{
    name={Aufgaben zur Konzentrationsberechnung, sowohl \% als auch Liter als auch Mol},
    description={todo},
    symbol={2015, 2018},
    type=frage
}

\newglossaryentry{xgal}{
    name={X-Gal},
    description={todo},
    symbol={2018},
    type=frage
}

\newglossaryentry{assemblierung}{
    name={Erkläre was eine Assemblierung ist},
    description={todo},
    symbol={2018},
    type=frage
}

\newglossaryentry{multiplesalignment}{
    name={Skizziere ein multiples Alignment},
    description={todo},
    symbol={2018},
    type=frage
}

\newglossaryentry{sequenzenunterscheiden}{
    name={Gegegeben zwei Sequenzen mit einer Insertion. Nenne zwei Strategien, um diese mit Primern auf dem Gel zu unterscheiden.},
    description={todo},
    symbol={2018},
    type=frage
}



\newglossaryentry{ptclvorteile}{
    name={Je zwei Vor- und Nachteile von PTCL?},
    description={todo},
    symbol={2018 (Vorlesung)},
    type=frage
}

\newglossaryentry{durchführungPTCL}{
    name={Durchführung PTCL},
    description={todo},
    symbol={2018 (Vorlesung)},
    type=frage
}

\newglossaryentry{HPTCLvsPTCL}{
    name={Unterschiede HPTCL und PTCL},
    description={todo},
    symbol={2018 (Vorlesung)},
    type=frage
}

\newglossaryentry{HPTCL}{
    name={Was ist HPTCL? / Wie funktioniert es?},
    description={todo},
    symbol={2015, 2018 (Vorlesung)},
    type=frage
}

\newglossaryentry{zwischentranskriptionundtranslation}{
    name={Schritte zwischen Translation und Transkription (4) mit zugehöriger RNA-Zwischenstufe (5) nennen},
    description={todo},
    symbol={2018 (Vorlesung)},
    type=frage
}

\newglossaryentry{repeatprobleme}{
    name={Bioinformatik: Pobleme bei Repeats},
    description={todo},
    symbol={2020 (Praktikum)},
    type=frage
}

\newglossaryentry{genefinding}{
    name={Methoden zum Gene finding},
    description={todo},
    symbol={2020 (Praktikum)},
    type=frage
}

\newglossaryentry{fusionunterschiede}{
    name={Unterschied transkriptionelle/translationale Fusion},
    description={todo},
    symbol={2020 (Praktikum)},
    type=frage
}

\newglossaryentry{deletionenfinden}{
    name={Deletionen mithilfe von Primerkombinationen finden},
    description={todo},
    symbol={2020 (Praktikum)},
    type=frage
}

\newglossaryentry{rtpcr}{
    name={RT-PCR Schritte},
    description={todo},
    symbol={2020 (Praktikum)},
    type=frage
}

\newglossaryentry{rekombinatenmarker}{
    name={Rekombinanten und Marker (Exp 4)},
    description={todo},
    symbol={2020 (Praktikum)},
    type=frage
}


\newglossaryentry{stoffwechselwegmutante}{
    name={Stoffwechselweg-Mutanten (Enzyme und TF) (Exp 5)},
    description={todo},
    symbol={2020 (Praktikum)},
    type=frage
}
\newglossaryentry{innateimunity}{
    name={Immunitätstypen, MAMPs und Rezeptor bei innate imunity},
    description={todo},
    symbol={2020 (Vorlesung)},
    type=frage
}

\newglossaryentry{codingfinden}{
    name={Nicht alle Bereiche mit Start und
    stop codons codieren für
    Proteine. Wie kann man
    bioinformatisch rausfinden
    welche es tun und welche nicht?},
    description={todo},
    symbol={2025 (Praktikum)},
    type=frage
}

\newglossaryentry{blastsmithwaterman}{
    name={Warum ist Blast so gut also
    worauf basiert der Algorithmus
    im Vergleich zu Smith-
    Waterman Algorithmus?},
    description={todo},
    symbol={2025 (Praktikum)},
    type=frage
}

\newglossaryentry{alignmentunterschiede}{
    name={Unterschied lokales und
    globales Alignment?},
    description={todo},
    symbol={2025 (Praktikum)},
    type=frage
}

\newglossaryentry{bioinf}{
    name={N50 Wert Bedeutung, hoher
    oder niedriger Wert besser?},
    description={todo},
    symbol={2002, 2025 (Praktikum)},
    type=frage
}

\newglossaryentry{rnahäufigkeit}{
    name={Welche Art von RNA liegt am meisten
    vor},
    description={todo nicht mRNA davon
    gibt’s nur 1 bis 4 Prozent)},
    symbol={2025 (Praktikum)},
    type=frage
}

\newglossaryentry{kreuzungauswertung}{
    name={Auswerten einer F2 Kreuzung
    von Pflanze A und B mit 3
    verschiedenen Markern. Wie
    viele Rekombinationsereignisse?},
    description={todo},
    symbol={2025 (Praktikum)},
    type=frage
}

\newglossaryentry{HPTLCauswertung}{
    name={Auswertung von HPTLC: welche
    Moleküle werden bei welchen
    Transkriptionsfaktoren
    synthetisiert oder nicht?},
    description={todo},
    symbol={2025 (Praktikum)},
    type=frage
}

\newglossaryentry{gusfärbung}{
    name={GUS Färbung skizzieren},
    description={todo},
    symbol={2025 (Praktikum)},
    type=frage
}

\newglossaryentry{gusreporter}{
    name={Reportergene der GUS-Färbung},
    description={todo},
    symbol={2018 (Praktikum)},
    type=frage
}

\newglossaryentry{gusschritte}{
    name={GUS im Mutanten feststellen - Schritte + Färbung erklären},
    description={todo},
    symbol={2015, 2018, 2020 (Praktikum)},
    type=frage
}

\newglossaryentry{guspromotor}{
    name={GUS Promotoren Fusion, Arten},
    description={todo},
    symbol={2020 (Praktikum)},
    type=frage
}

\newglossaryentry{genexpressionsanalyse}{
    name={Methoden und deren Ablauf},
    description={todo},
    symbol={2025 (Praktikum)},
    type=frage
}

\newglossaryentry{rechenaufgabenII}{
    name={Rechenaufgaben},
    description={todo Rechenaufgaben mit den
    Formeln (c1 * v1 = c2 * v2) und
    m= n v M, mM zu M umrechnen
    dran denken und logisch denken
    er baut manchmal Fallen ein
    einfach genau schauen!!
    (Insgesamt 3 verschiedene
    Aufgaben)},
    symbol={2025 (Praktikum)},
    type=frage
}

\newglossaryentry{rechenaufgaben}{
    name={die beiden Formeln aus dem Praktikum anwenden},
    description={todo},
    symbol={2020 (Praktikum)},
    type=frage
}

\newglossaryentry{aktivtransposon}{
    name={Was haben aktive Transposons
    was nicht aktive nicht haben},
    description={todo},
    symbol={2025 (Vorlesung)},
    type=frage
}

\newglossaryentry{autonomietransposon}{
    name={Unterschied und Gemeinsamkeiten autonome und nicht
    autonome Transposons},
    description={todo},
    symbol={2020, 2025 (Vorlesung)},
    type=frage
}

\newglossaryentry{eukprounterschied}{
    name={Unterschiede zwischen Eukaryoten und Prokaryoten},
    description={todo Introns
    Vorhandensein, poly/
    monocystronisch, RNA-
    Polymerasen, Operons usw},
    symbol={2020, 2025 (Vorlesung)},
    type=frage
}

\newglossaryentry{klregrna}{
    name={(Drei) kleine regulatorische RNAs nennen und deren Funktionen ausschreiben},
    description={todo iRNA, miRNA, piRNA usw},
    symbol={2020, 2025 (Vorlesung)},
    type=frage
}

\newglossaryentry{MAMPPRR}{
    name={ MAMP und PRR
    ausschreiben und deren Funktion nennen},
    description={todo},
    symbol={2025 (Vorlesung)},
    type=frage
}


\newglossaryentry{abbgen}{
    name={Abbildung eines Gens beschriften},
    description={todo},
    symbol={2020, 2025 (Vorlesung)},
    type=frage
}

\newglossaryentry{abbmrna}{
    name={Abbildung einer mRNA beschriften},
    description={todo},
    symbol={2020, 2025 (Vorlesung)},
    type=frage
}

\newglossaryentry{opineAgrobakterium}{
    name={Wozu sind Opine im Agrobakterium da},
    description={todo},
    symbol={2025 (Vorlesung)},
    type=frage
}

\newglossaryentry{transformationsubstanzen}{
    name={Substanzen von transformierten Zellen (Opine?), Funktion usw.},
    description={todo},
    symbol={2020 (Vorlesung)},
    type=frage
}

\newglossaryentry{einbuchstabencode}{
    name={Aminosäure-Code Zuordnung (bspw. KALTE (2025), FLEISS (2015, 2018))},
    description={todo},
    symbol={2015, 2018, 2020, 2025 (Vorlesung)},
    type=frage
}


\newglossaryentry{lückentextspleißosom}{
    name={Lückentext zu nuclearem
    spleißen: introns, spleißson,
    Aufbau Spleißosom,},
    description={todo},
    symbol={2025 (Vorlesung)},
    type=frage
}



%% eingebaut:

\newglossaryentry{transportsignaleimport}{
    name={4 Transportsignale bei: Import
    ER, Retention Er, Mitochondrien
    und Plastom, Lysosom, nennen
    und ausschreiben},
    description={todo},
    symbol={2025 (Vorlesung)},
    type=frage
}

\newglossaryentry{transportsignaleimportII}{
    name={ Proteinimport Signale - Kern, Retention im ER, Lysosome, Mitochondrien},
    description={todo},
    symbol={2020 (Vorlesung)},
    type=frage
}

\newglossaryentry{proteinsortingsignale}{
    name={Die drei Signale vom Protein Sorting mit Beschreibung},
    description={todo},
    symbol={2015, 2018 (Vorlesung)},
    type=frage
}

\newglossaryentry{signaleproteinsorting}{
    name={Signale des Protein-Sorting mit Beschreibung},
    description={todo},
    symbol={2018 (Vorlesung)},
    type=frage
}